% !TEX root = ../main.tex
% Бүлэг 8 — Дүгнэлт

\chapter{Дүгнэлт}
\label{Conclusion}
\pagecolor{white}

\section*{Ерөнхий дүгнэлт}

Энэхүү судалгааны ажлын хүрээнд рентген зурагт суурилан ясны хугарлыг ангилах, сегментлэх, хугарлын зэргийг үнэлэх болон эдгэрэлтийн төлөвийг таамаглах боломжийг гүн сургалтын аргаар судалж, олон үе шаттай хиймэл оюуны pipeline боловсрууллаа. Судалгааны явцад нэгдсэн олон даалгаварт архитектур, дараалсан сургалтын арга болон модульчлагдсан pipeline гэсэн гурван өөр хандлагыг туршин харьцуулсан бөгөөд эцсийн байдлаар бие даасан модулиудаас бүрдэх pipeline хамгийн тогтвортой, найдвартай үр дүн үзүүлсэн.

Туршилтын үр дүнгээр ConvNeXt-Base ангилагч нь 3-fold cross-validation дээр AUC-ROC = 0.9887, F1 = 0.959 үзүүлэлтэд хүрч, судалгаанд ашигласан бусад ангилагч загваруудаас өндөр гүйцэтгэл үзүүлэв. Мөн ResNet50 U-Net сегментчлэлийн загвар нь тест олонлог дээр Dice = 0.604 үр дүн гаргаж, анхны нэгдсэн архитектурын Dice = 0.0498 үзүүлэлтээс мэдэгдэхүйц сайжирсан.

Cургалтад огт ороогүй \textbf{hold-out тестийн санд} ангиллын Test AUC = 0.9703 (generalization gap = 0.018), F1 = 0.804; pipeline-conditional сегментчлэлийн Dice = 0.477 (TP=43) гарсан нь загвар бэрхшээлгүй ерөнхийлж байгааг харуулна. Үүнээс гадна, Stanford ML Group-ийн MURA-v1.1 өгөгдөл дээр \textbf{гадаад (external) валидаци} хийсэн 3{,}197 зурганд Test AUC = 0.8323 (FOREARM 0.93, HUMERUS 0.93, WRIST 0.87) авч pipeline-ын domain-аас гадуурх өгөгдөлд ажиллах чадавхыг баталгаажуулсан. Энэ нь хугарлын ангилал болон сегментчлэлийн даалгавруудыг нэг кодлогч дээр зэрэг сургах үед task interference үүсэж болох бөгөөд тодорхой нөхцөлд модульчлагдсан архитектур илүү үр ашигтай байж болохыг харуулж байна.

Сегментчлэлийн үр дүнд үндэслэн гарган авсан морфологийн шинжүүдийг ашиглан хугарлын зэргийг proof-of-concept түвшинд үнэлэх боломжтойг судалгааны үр дүн харуулсан. Түүнчлэн эдгэрэлтийн хугацаатай холбоотой хүчин зүйлсийг ашиглан эдгэрэлтийн төлөвийг таамаглах proof-of-concept туршилт хийсэн нь хугарлын автомат шинжилгээг зөвхөн илрүүлэлтээр хязгаарлахгүйгээр дараагийн шатны клиник мэдээлэлтэй холбох боломж байгааг харуулсан.

Судалгааны ажлын хүрээнд FracAtlas, GRAZPEDWRI-DX, FracAtlas-COCO, YOLOBoneFracture болон BoneFractureCVProject зэрэг олон нийтэд нээлттэй өгөгдлийн сангуудыг нэгтгэн ашиглаж, ясны хугарлын ангилал, сегментчлэл болон дараагийн шатны үнэлгээний асуудлуудыг цогцоор нь судалсан. Мөн Grad-CAM болон SHAP тайлбарлах аргуудыг ашигласнаар загварын шийдвэр гаргалтыг илүү ойлгомжтой тайлбарлах боломжийг харуулсан.

Гэсэн хэдий ч энэхүү судалгаа нь тодорхой хязгаарлалтуудтай. Ашигласан өгөгдлийн сангууд нь Монгол хүн амын өгөгдлийг агуулаагүй бөгөөд сегментчлэлийн сургалтад ашигласан зарим маск автоматаар үүсгэсэн псевдо шошго байсан. Хугарлын зэрэг үнэлгээ нь бодит клиникийн хугарлын зэргийн шошго бус морфологийн шинжүүдэд тулгуурласан proof-of-concept үнэлгээ бөгөөд эдгэрэлтийн модуль нь бодит follow-up өгөгдөл бус синтетик өгөгдөл дээр суурилсан proof-of-concept туршилт юм. Иймээс эдгээр үр дүнг клиникийн орчинд шууд хэрэглэхээс өмнө бодит эмнэлгийн өгөгдөл дээр дахин баталгаажуулах шаардлагатай.

Ерөнхийд нь дүгнэвэл энэхүү ажил нь ясны хугарлыг ангилах, сегментлэх, хугарлын зэргийг үнэлэх болон эдгэрэлтийн төлөвийг судлах хиймэл оюуны цогц хандлагыг боловсруулж, модульчлагдсан архитектурын давуу талыг туршилтаар харуулсан судалгааны ажил болсон. Цаашид бодит эмнэлгийн өгөгдөл, клиникийн баталгаажуулалт болон олон төрлийн дүрслэлийн мэдээлэлтэй хослуулан хөгжүүлснээр ясны хугарлын оношилгооны чанар, хүртээмжийг сайжруулахад хувь нэмэр оруулах боломжтой гэж үзэж байна.
